\chapter{非线性数据结构：哈希表、堆、树、图和并查集}

线性结构擅长处理顺序，非线性结构擅长处理关系和查询。哈希表回答“这个状态以前出现过吗”；堆回答“当前最极端的候选是谁”；树表达层次、有序性和递归子问题；图表达任意状态之间的连接；并查集维护“这些元素现在是否属于同一组”。这些结构不是刷题技巧的附属品，而是算法优化的来源。很多题从 \complexity{O(n^2)} 变成 \complexity{O(n)}，并不是因为循环少写了一层，而是因为某个昂贵查询被换成了合适的数据结构查询。

哈希表是保存历史信息的核心结构。\code{std::unordered_map} 和 \code{std::unordered_set} 的平均查询、插入、删除复杂度是 \complexity{O(1)}，适合把“回头查找”变成“即时查询”。两数之和的暴力方法枚举两个下标，慢在每个数都要向后找补数；哈希表版本边扫描边记录已经见过的值，当前数只需要查询一次补数。这个优化成立的前提是题目只要求找到一对下标，且补数是否出现过可以由历史集合回答。

\begin{lstlisting}[language=C++,caption={两数之和：哈希表把补数查询变成平均常数时间}]
std::vector<int> two_sum(const std::vector<int>& nums, int target)
{
    std::unordered_map<int, int> index_by_value;
    index_by_value.reserve(nums.size());

    for (int i = 0; i < static_cast<int>(nums.size()); ++i) {
        const int complement = target - nums[i];
        const auto found = index_by_value.find(complement);
        if (found != index_by_value.end()) {
            return {found->second, i};
        }
        index_by_value.emplace(nums[i], i);
    }

    return {};
}
\end{lstlisting}

代码里先查再插入，是为了避免同一个元素被用两次。例如目标是 6，当前数是 3，如果先插入再查，可能错误地返回同一个下标。这里用 \code{find} 而不是 \code{operator[]}，因为 \code{operator[]} 在 key 不存在时会插入默认值。只想判断存在性时，C++20 可以用 \code{contains}；需要同时取值时，用 \code{find}，这样只查一次。刷题中还应该习惯 \code{reserve}，它能减少 rehash 次数。rehash 会重新组织桶，导致迭代器失效；如果保存了哈希表元素迭代器，扩容后不能继续使用。

哈希表不仅能存值到下标，也能存频次、前缀和、状态编码和访问标记。前缀和加哈希表是很重要的一类优化。和为 K 的子数组如果暴力枚举左右端点，需要 \complexity{O(n^2)}；如果知道当前位置前缀和是 \code{current}，那么以当前位置结尾、和为 \code{k} 的子数组数量，等于以前出现过多少次 \code{current - k}。哈希表保存的是“某个前缀和出现次数”，查询的是“能和当前前缀配成目标的历史前缀”。

\begin{lstlisting}[language=C++,caption={和为 K 的子数组：前缀和频次}]
int subarray_sum_equals_k(const std::vector<int>& nums, int k)
{
    std::unordered_map<int, int> prefix_count;
    prefix_count.reserve(nums.size() + 1);
    prefix_count.emplace(0, 1);

    int prefix = 0;
    int answer = 0;
    for (const int x : nums) {
        prefix += x;
        const auto found = prefix_count.find(prefix - k);
        if (found != prefix_count.end()) {
            answer += found->second;
        }
        ++prefix_count[prefix];
    }

    return answer;
}
\end{lstlisting}

这个算法容易错在初始化。空前缀和 0 必须先出现一次，否则从下标 0 开始的合法子数组会漏掉。它也解释了为什么有负数时不能随便用滑动窗口：窗口和不再随右端单调增加，左端移动也不保证更接近目标；前缀和加哈希表不依赖正数单调性，所以适用范围更广。

哈希表的限制同样要清楚。第一，它不维护顺序，如果题目要求按排序顺序取最小、最大或区间，就该考虑 \code{std::map}、\code{std::set}、堆、树状数组或线段树。第二，平均 \complexity{O(1)} 不等于绝对 \complexity{O(1)}，糟糕哈希或恶意数据可能退化。第三，自定义 key 需要提供相等比较和哈希函数。面试里一般不会刻意卡哈希攻击，但你要知道 \code{unordered_map<pair<int,int>, int>} 不能直接在所有标准库环境下工作，需要自定义哈希。

\begin{lstlisting}[language=C++,caption={为 pair 坐标提供哈希函数}]
struct PairHash {
    std::size_t operator()(const std::pair<int, int>& point) const noexcept
    {
        constexpr unsigned COORDINATE_BITS = 32U;
        const auto first = static_cast<std::uint64_t>(
            static_cast<std::uint32_t>(point.first));
        const auto second = static_cast<std::uint64_t>(
            static_cast<std::uint32_t>(point.second));
        return std::hash<std::uint64_t>{}((first << COORDINATE_BITS) ^ second);
    }
};

std::unordered_set<std::pair<int, int>, PairHash> visited;
\end{lstlisting}

堆解决的是动态极值问题。C++ 里通常使用 \code{std::priority_queue}，默认是大顶堆；如果要小顶堆，需要显式给出容器和比较器。堆的插入和弹出堆顶都是 \complexity{O(\log n)}，查看堆顶是 \complexity{O(1)}。它不适合频繁删除任意元素，也不适合遍历出有序结果而不破坏结构。换句话说，堆只承诺“我能很快告诉你当前最大或最小是谁”。

\begin{lstlisting}[language=C++,caption={priority\_queue 的大顶堆和小顶堆}]
std::priority_queue<int> max_heap;

std::priority_queue<int,
                    std::vector<int>,
                    std::greater<int>> min_heap;

max_heap.push(3);
max_heap.push(10);
const int largest = max_heap.top();
max_heap.pop();
\end{lstlisting}

Top K 是堆最经典的用法。求数组中第 K 大元素，直觉做法是排序，复杂度 \complexity{O(n\log n)}。如果只关心前 K 大，不需要完整排序，可以维护一个大小为 K 的小顶堆。堆顶是当前前 K 大中最小的元素；新元素如果不超过堆顶，就不可能进入前 K；如果超过堆顶，就弹出堆顶并插入新元素。最终堆顶就是第 K 大。

\begin{lstlisting}[language=C++,caption={第 K 大元素：大小为 K 的小顶堆}]
int find_kth_largest(const std::vector<int>& nums, int k)
{
    std::priority_queue<int,
                        std::vector<int>,
                        std::greater<int>> heap;

    for (const int x : nums) {
        if (static_cast<int>(heap.size()) < k) {
            heap.push(x);
        } else if (x > heap.top()) {
            heap.pop();
            heap.push(x);
        }
    }

    return heap.top();
}
\end{lstlisting}

这个方法的复杂度是 \complexity{O(n\log k)}，适合 k 远小于 n 的场景。还有一种平均线性复杂度的快速选择算法，后文会讲。面试表达时要说明选择堆的理由：不需要完整有序，只需要维护 K 个候选的边界。类似地，合并 K 个升序链表用小顶堆保存每条链表当前头节点；任务调度用堆保存当前最早结束或最高优先级任务；数据流中位数用一个大顶堆和一个小顶堆分别维护较小一半和较大一半。

树把数据组织成层次关系。二叉树题看似在考递归，实际在考状态如何沿边传播。状态可以从父节点传给子节点，例如验证二叉搜索树时传递取值上下界；也可以从子节点汇总到父节点，例如求树的高度、直径、最大路径和；还可以按层从左到右传播，例如层序遍历。只背前序、中序、后序三个词是不够的，必须知道遍历顺序服务于哪类状态。

\begin{lstlisting}[language=C++,caption={二叉搜索树验证：显式栈中序遍历}]
bool is_valid_bst(TreeNode* root)
{
    std::vector<TreeNode*> stack;
    TreeNode* current = root;
    std::optional<long long> previous;

    while (current != nullptr || !stack.empty()) {
        while (current != nullptr) {
            stack.push_back(current);
            current = current->left;
        }

        current = stack.back();
        stack.pop_back();

        if (previous.has_value() && current->value <= previous.value()) {
            return false;
        }
        previous = current->value;
        current = current->right;
    }

    return true;
}
\end{lstlisting}

二叉搜索树的关键性质是中序遍历严格递增。验证 BST 的暴力想法可能是对每个节点扫描整棵左子树和右子树，判断左边都小、右边都大，这会重复访问大量节点。中序遍历把全局约束转成相邻访问值的递增关系，每个节点只处理一次。这里用显式栈而不是递归，是为了避免极端深树导致 C++ 调用栈过深，同时也让“遍历状态”更清楚。

树题经常需要区分“节点为空的返回值”和“答案的初始值”。求最大深度时，空节点深度是 0；求最大路径和时，答案不能初始化为 0，因为所有节点可能都是负数；求最近公共祖先时，返回值表示“当前子树是否找到了目标节点或祖先”。这些不是小细节，而是树形 DP 的语义。每个返回值都应该能用一句话解释：它代表当前子树对父节点贡献什么信息。

图是最通用的关系结构。只要有状态和一步转移，就可以建模成图。图题最重要的不是 DFS 或 BFS 写法，而是建模四问：节点是什么，边是什么，访问状态是什么，队列或栈里保存什么。岛屿数量里，节点是陆地格子，边是上下左右相邻；课程表里，节点是课程，有向边表示先修依赖；单词接龙里，节点是单词，边表示一次字符变化；开锁问题里，节点是四位密码状态，边是转动一位数字。

图的存储方式要跟输入规模匹配。稠密图可以用邻接矩阵，判断两点是否有边是 \complexity{O(1)}，但空间是 \complexity{O(n^2)}。大多数面试题是稀疏图，更适合邻接表，空间是 \complexity{O(n + m)}。如果节点编号连续，用 \code{vector<vector<int>>}；如果节点是字符串或稀疏编号，用哈希表映射到编号，或者直接用 \code{unordered_map<string, vector<string>>}。

\begin{lstlisting}[language=C++,caption={课程表：邻接表和入度数组}]
bool can_finish(int num_courses,
                const std::vector<std::vector<int>>& prerequisites)
{
    std::vector<std::vector<int>> graph(num_courses);
    std::vector<int> indegree(num_courses, 0);

    for (const auto& edge : prerequisites) {
        const int course = edge[0];
        const int prerequisite = edge[1];
        graph[prerequisite].push_back(course);
        ++indegree[course];
    }

    std::queue<int> queue;
    for (int course = 0; course < num_courses; ++course) {
        if (indegree[course] == 0) {
            queue.push(course);
        }
    }

    int learned = 0;
    while (!queue.empty()) {
        const int course = queue.front();
        queue.pop();
        ++learned;

        for (const int next : graph[course]) {
            --indegree[next];
            if (indegree[next] == 0) {
                queue.push(next);
            }
        }
    }

    return learned == num_courses;
}
\end{lstlisting}

拓扑排序的暴力直觉是反复扫描所有课程，找当前没有前置依赖的课。这样每学一门课都可能重新扫描所有边。入度数组把“还有多少依赖没被满足”变成可增量更新的状态；队列保存当前入度为 0 的课程。每条边只在它的起点被移除时处理一次，复杂度 \complexity{O(n + m)}。如果最后学完的课程数少于总数，说明剩余节点互相依赖，图中有环。

DFS 和 BFS 都可以遍历图，但语义不同。BFS 适合无权最短步数和层次扩散，因为它按距离递增访问状态；DFS 适合连通块标记、路径搜索、拓扑状态检测和回溯枚举，因为它沿一条路径深入。C++ 里递归 DFS 简洁，但深度过大可能栈溢出。本书讲递归思想，但实现上会经常给出显式栈版本，让遍历过程、访问标记和状态回退更加可控。

\begin{lstlisting}[language=C++,caption={岛屿数量：显式栈标记连通块}]
int num_islands(std::vector<std::vector<char>>& grid)
{
    if (grid.empty() || grid[0].empty()) {
        return 0;
    }

    constexpr char LAND = '1';
    constexpr char WATER = '0';
    constexpr std::array<std::pair<int, int>, 4> DIRECTIONS{
        std::pair{1, 0}, std::pair{-1, 0},
        std::pair{0, 1}, std::pair{0, -1}
    };

    const int rows = static_cast<int>(grid.size());
    const int cols = static_cast<int>(grid[0].size());
    int islands = 0;

    for (int row = 0; row < rows; ++row) {
        for (int col = 0; col < cols; ++col) {
            if (grid[row][col] != LAND) {
                continue;
            }

            ++islands;
            std::vector<std::pair<int, int>> stack{{row, col}};
            grid[row][col] = WATER;

            while (!stack.empty()) {
                const auto [r, c] = stack.back();
                stack.pop_back();

                for (const auto [dr, dc] : DIRECTIONS) {
                    const int nr = r + dr;
                    const int nc = c + dc;
                    if (nr < 0 || nr >= rows || nc < 0 || nc >= cols ||
                        grid[nr][nc] != LAND) {
                        continue;
                    }
                    grid[nr][nc] = WATER;
                    stack.emplace_back(nr, nc);
                }
            }
        }
    }

    return islands;
}
\end{lstlisting}

这段代码把访问过的陆地直接改成水，省掉单独的 \code{visited} 数组。这样做是否合适取决于题目是否允许修改输入。如果不允许，就创建同样大小的布尔矩阵。图题里访问标记必须在入队或入栈时就设置，而不是弹出时才设置；否则同一个节点可能被多个邻居重复加入，轻则浪费时间，重则导致状态爆炸。

并查集维护动态连通性。它有两个核心操作：\code{find} 找到一个元素所在集合的代表，\code{unite} 合并两个集合。路径压缩让查找时顺手把沿途节点直接挂到根上，按大小或按秩合并让树保持较浅。两者结合后，单次操作的均摊复杂度接近常数。它适合回答“是否在同一组”“合并后还有几个连通块”，不适合回答两点之间的具体路径或最短距离。

\begin{lstlisting}[language=C++,caption={并查集：路径压缩与按大小合并}]
class DisjointSetUnion {
public:
    explicit DisjointSetUnion(int n)
        : parent_(n), size_(n, 1)
    {
        std::iota(this->parent_.begin(), this->parent_.end(), 0);
    }

    int find(int x)
    {
        int root = x;
        while (this->parent_[root] != root) {
            root = this->parent_[root];
        }
        while (this->parent_[x] != x) {
            const int next = this->parent_[x];
            this->parent_[x] = root;
            x = next;
        }
        return root;
    }

    bool unite(int a, int b)
    {
        int root_a = find(a);
        int root_b = find(b);
        if (root_a == root_b) {
            return false;
        }
        if (this->size_[root_a] < this->size_[root_b]) {
            std::swap(root_a, root_b);
        }
        this->parent_[root_b] = root_a;
        this->size_[root_a] += this->size_[root_b];
        return true;
    }

private:
    std::vector<int> parent_;
    std::vector<int> size_;
};
\end{lstlisting}

冗余连接是并查集的标准例题。给一棵树额外加一条边，依次扫描边。如果一条边连接的两个点已经在同一集合里，那么再加入它就会形成环，这条边就是冗余边。暴力方法可以每次加边后做一次 DFS 检查连通性，复杂度很高；并查集把“两个点是否已经连通”维护成近乎常数时间查询。

账户合并也能体现并查集的建模能力。邮箱是连接账户的证据：两个账户只要共享一个邮箱，就属于同一个人。可以把账户编号作为并查集元素，用哈希表记录邮箱第一次出现在哪个账户；再次看到同一邮箱时，把两个账户合并。最后按根节点收集邮箱并排序。这里哈希表负责从字符串邮箱找到账户编号，并查集负责维护账户之间的连通分量。一个题常常需要多个数据结构协作，这也是“数据结构和算法不分家”的真实含义。

\begin{principlebox}{数据结构和算法不分家}
优化通常来自查询模型的改变。哈希表把历史查找变快，堆把动态极值变快，树把层次状态和有序性质显式化，图把状态转移统一成边，并查集把动态连通性维护变快。刷题复盘时必须同时写出算法动作和数据结构职责：保存什么，支持什么查询，删除什么过期状态，为什么不会漏答案。
\end{principlebox}

本章建议配套练习：哈希表做 1、49、128、560、146；堆做 215、347、23、295、621；树做 94、98、102、104、105、124、236、230；图做 200、207、210、417、994、127、752；并查集做 547、684、721、990、1319。每道题至少复盘一次建模：如果不用这个结构，暴力瓶颈在哪里；用了这个结构，哪一步查询或合并被降下来了；结构本身有什么边界和失效条件。

\section{非线性结构的教材级展开}

非线性结构解决的是“关系查询”而不只是“顺序扫描”。哈希表把对象映射到桶，堆把候选组织成局部有序的完全二叉树，平衡树维护全局有序关系，Trie 把字符串前缀共享起来，树状数组和线段树把区间信息按层聚合，图把任意状态连接成网络，并查集把动态连通关系压缩成集合代表。经典教材会把这些结构讲得很细，因为它们不仅出现在面试题里，也出现在编译器符号表、数据库索引、操作系统调度、网络路由、文件系统目录和缓存系统中。

学习非线性结构要警惕一个误区：只记复杂度表。复杂度表只告诉你某个操作大概多快，却没有告诉你结构为什么能这么快、在哪些操作上不擅长、变体题需要保存哪些额外信息。比如哈希表平均查询常数，但无法找前驱后继；堆能快速取极值，但不能快速删除任意元素；平衡树能维护有序，但常数比数组扫描大；并查集能判断同组，却不能告诉你最短路径。高质量题解必须把能力边界讲清楚。

\section{哈希表：桶、冲突、负载因子和历史状态}

哈希表的核心思想是把 key 通过哈希函数映射到桶。理想情况下，key 均匀分布到桶中，每个桶里元素很少，查询接近常数时间。实际实现需要处理冲突：常见方式有链地址法和开放寻址。C++ 的 \code{std::unordered_map} 抽象上承诺平均常数时间，具体实现通常围绕桶数组、节点和负载因子展开。当元素数量相对桶数量过高时，容器会 rehash，重新分配桶并把元素放到新桶里。

刷题中，哈希表最常见的职责是保存历史状态。两数之和保存历史值到下标，前缀和题保存历史前缀出现次数，最长连续序列保存所有值的存在性，快乐数保存已经出现过的状态，克隆图保存原节点到新节点的映射。虽然都叫哈希表，保存内容完全不同：有时保存布尔存在，有时保存次数，有时保存最早位置，有时保存对象指针。题目目标决定哈希表的 value 类型。

同一题型的变体常常只改哈希表语义。和为 K 的子数组问数量，哈希表保存前缀和频次；问最长长度，保存前缀和最早出现位置；问是否存在长度至少二的倍数子数组，保存余数最早出现位置；问最短且允许负数，则普通哈希不够，需要前缀和加单调队列。复盘时不能只写“用哈希表”，必须写“哈希表保存什么历史信息”。

\begin{lstlisting}[language=C++,caption={最长和为 K 的子数组：保存最早前缀位置}]
int longest_subarray_sum_k(const std::vector<int>& nums, int target)
{
    std::unordered_map<std::int64_t, int> first_index;
    first_index.reserve(nums.size() + 1);
    first_index.emplace(std::int64_t{0}, -1);

    std::int64_t prefix = 0;
    int answer = 0;
    for (int i = 0; i < static_cast<int>(nums.size()); ++i) {
        prefix += nums[i];
        const auto found = first_index.find(prefix - target);
        if (found != first_index.end()) {
            answer = std::max(answer, i - found->second);
        }
        first_index.emplace(prefix, i);
    }
    return answer;
}
\end{lstlisting}

这里使用 \code{emplace} 而不是直接赋值，是为了保留某个前缀和第一次出现的位置。要最长区间，左端越早越好；如果每次更新为最新位置，就会把潜在更长答案破坏掉。这说明同样是前缀和哈希，计数题、最长题、最短题的 value 语义完全不同。

哈希表的实际应用远超刷题。编译器符号表用哈希表从名字找到声明信息；解释器运行时环境用哈希表表示对象属性；缓存系统用哈希表定位缓存项；数据库执行层用哈希表做 hash join；操作系统内核用哈希表或类似结构快速定位对象。应用越工程化，越要关心哈希函数、内存分配、并发访问和迭代器失效。刷题可以先掌握模型，进阶时必须知道平均常数背后的条件。

\section{堆和优先队列：动态极值、懒删除和多路归并}

二叉堆通常用数组表示完全二叉树。下标 i 的左右孩子在 \code{2*i+1} 和 \code{2*i+2}，父节点在 \code{(i-1)/2}。大顶堆要求父节点不小于孩子，小顶堆要求父节点不大于孩子。堆不是完全有序结构，只保证堆顶是全局极值。插入时向上调整，弹出堆顶时把末尾元素放到根再向下调整，复杂度都是 \complexity{O(\log n)}。

优先队列适合“候选集合动态变化，但每次只需要当前最优候选”的问题。Top K、合并 K 个有序链表、数据流中位数、任务调度、Dijkstra、最小区间覆盖 K 个列表、IPO 项目选择都属于这一类。堆不适合的场景是频繁查找或删除任意元素，因为普通 \code{priority_queue} 只能访问堆顶。

懒删除是堆的重要变体。滑动窗口中位数需要删除窗口左端过期元素，但堆无法直接删除任意元素。常见做法是用哈希表记录某个值应该被删除的次数，真正访问堆顶前，不断弹出已标记删除的堆顶。这种策略把任意删除延迟到元素成为堆顶时执行，保持堆接口简单。Dijkstra 中的过期距离也是同样思想：不更新堆中旧节点，而是压入新距离，弹出时检查是否过期。

\begin{lstlisting}[language=C++,caption={堆中旧状态懒跳过：Dijkstra 的核心形态}]
std::vector<std::int64_t> dijkstra(
    const std::vector<std::vector<std::pair<int, int>>>& graph,
    int source)
{
    constexpr std::int64_t INF = 4'000'000'000'000'000'000;
    using State = std::pair<std::int64_t, int>;

    std::vector<std::int64_t> distance(graph.size(), INF);
    std::priority_queue<State, std::vector<State>, std::greater<State>> heap;
    distance[static_cast<std::size_t>(source)] = 0;
    heap.emplace(0, source);

    while (!heap.empty()) {
        const auto [current_distance, node] = heap.top();
        heap.pop();
        if (current_distance != distance[static_cast<std::size_t>(node)]) {
            continue;
        }

        for (const auto& [next, weight] : graph[static_cast<std::size_t>(node)]) {
            const std::int64_t candidate = current_distance + weight;
            if (candidate >= distance[static_cast<std::size_t>(next)]) {
                continue;
            }
            distance[static_cast<std::size_t>(next)] = candidate;
            heap.emplace(candidate, next);
        }
    }
    return distance;
}
\end{lstlisting}

堆的变种包括二叉堆、d 叉堆、可并堆、斐波那契堆、索引堆和双堆。面试中通常只需要二叉堆和双堆。双堆维护中位数时，一个大顶堆保存较小一半，一个小顶堆保存较大一半；不变量是大顶堆所有值不大于小顶堆所有值，两个堆大小差不超过一。很多错误实现只维护大小不变量，忘了顺序不变量，最终中位数会错。

\section{有序结构：平衡树、前驱后继和区间维护}

哈希表不能回答顺序问题。若题目需要找小于等于某值的最大元素、大于等于某值的最小元素、动态维护区间、按时间顺序查询，通常要考虑有序结构。C++ 的 \code{std::map} 和 \code{std::set} 通常由红黑树实现，插入、删除、查询和 \code{lower_bound} 都是 \complexity{O(\log n)}。它们常数比哈希表大，但能提供顺序语义。

有序结构典型应用包括日程安排、区间合并、滑动窗口内有序统计、在线前驱后继、时间戳键值存储、离线扫描线。比如 My Calendar 类问题中，要插入一个新区间，需要找到第一个起点不小于新区间起点的区间，再检查它和前一个区间是否冲突。哈希表无法回答“相邻区间是谁”，有序 map 正好能回答。

\begin{lstlisting}[language=C++,caption={有序 map 维护互不重叠区间}]
class IntervalSet {
public:
    bool add(int left, int right)
    {
        auto next = this->ranges_.lower_bound(left);
        if (next != this->ranges_.end() && next->first < right) {
            return false;
        }
        if (next != this->ranges_.begin()) {
            auto previous = std::prev(next);
            if (previous->second > left) {
                return false;
            }
        }
        this->ranges_.emplace(left, right);
        return true;
    }

private:
    std::map<int, int> ranges_;
};
\end{lstlisting}

这段代码体现了有序结构的独特能力：通过 \code{lower_bound} 直接定位可能冲突的相邻区间，而不是扫描所有区间。变体包括闭区间和半开区间的端点差异、允许重叠次数最多为二、插入后合并相邻区间、查询某点被多少区间覆盖。每个变体都要先定义端点语义，再决定比较条件里的等号。

\section{Trie：前缀共享和字符串集合}

Trie 把字符串集合组织成字符边组成的树。根到某节点的路径表示一个前缀，节点上的标记表示是否有单词在这里结束。它适合前缀查询、多模式搜索、自动补全、字典树剪枝和按位 Trie 最大异或。相比哈希表，Trie 的优势不是单个完整字符串查询更快，而是前缀共享和前缀剪枝。

实现 Trie 时要根据字符集选择孩子结构。小写英文可以用 \code{array<int, 26>} 或指针数组；字符集大时可以用哈希表保存孩子；若内存敏感，可以用压缩 Trie 或把节点存进 \code{vector} 里用整数下标引用。刷题中，单词搜索 II 是 Trie 的代表：如果对每个单词单独 DFS，会重复遍历棋盘；Trie 合并所有单词前缀后，DFS 一旦发现当前前缀不存在，就可以立即剪枝。

\begin{lstlisting}[language=C++,caption={Trie 节点用数组孩子表达小写字母前提}]
class Trie {
public:
    void insert(const std::string& word)
    {
        int node = 0;
        for (const char ch : word) {
            const int index = ch - 'a';
            if (this->nodes_[static_cast<std::size_t>(node)].next[index] == -1) {
                this->nodes_[static_cast<std::size_t>(node)].next[index] =
                    static_cast<int>(this->nodes_.size());
                this->nodes_.push_back(Node{});
            }
            node = this->nodes_[static_cast<std::size_t>(node)].next[index];
        }
        this->nodes_[static_cast<std::size_t>(node)].terminal = true;
    }

    bool starts_with(const std::string& prefix) const
    {
        int node = 0;
        for (const char ch : prefix) {
            const int index = ch - 'a';
            node = this->nodes_[static_cast<std::size_t>(node)].next[index];
            if (node == -1) {
                return false;
            }
        }
        return true;
    }

private:
    struct Node {
        std::array<int, 26> next{};
        bool terminal = false;

        Node()
        {
            this->next.fill(-1);
        }
    };

    std::vector<Node> nodes_{Node{}};
};
\end{lstlisting}

这里使用整数下标而不是裸指针，可以避免手动释放节点，也让节点存储更紧凑。代价是 \code{vector} 扩容时如果外部保存了节点引用会失效，所以代码只保存下标。字符集假设同样必须说明：输入如果不是小写英文，\code{ch - 'a'} 就不安全。

\section{树状数组和线段树：区间聚合和动态更新}

前缀和适合静态数组区间求和，但如果数组中间会更新，普通前缀和每次更新都要影响后面所有位置。树状数组和线段树解决的是“动态更新 + 区间查询”。树状数组适合可逆前缀聚合，代码短、常数小；线段树更通用，可以维护区间最大、最小、和、懒标记、复杂合并信息。

树状数组把每个位置负责的一段长度设置为最低位一的大小。更新一个点时，沿着 \code{index += lowbit(index)} 更新所有包含它的树节点；查询前缀时，沿着 \code{index -= lowbit(index)} 汇总若干不重叠块。它常用于逆序对、动态前缀和、离散化后的频次排名。

\begin{lstlisting}[language=C++,caption={树状数组：点更新和前缀查询}]
class FenwickTree {
public:
    explicit FenwickTree(int n) : tree_(static_cast<std::size_t>(n + 1), 0) {}

    void add(int index, std::int64_t delta)
    {
        for (int i = index + 1; i < static_cast<int>(this->tree_.size());
             i += i & -i) {
            this->tree_[static_cast<std::size_t>(i)] += delta;
        }
    }

    std::int64_t prefix_sum(int count) const
    {
        std::int64_t result = 0;
        for (int i = count; i > 0; i -= i & -i) {
            result += this->tree_[static_cast<std::size_t>(i)];
        }
        return result;
    }

private:
    std::vector<std::int64_t> tree_;
};
\end{lstlisting}

这段代码使用外部零基下标，内部一基下标。外部 \code{add(index, delta)} 更新第 index 个元素，\code{prefix_sum(count)} 查询前 count 个元素之和。把接口设计成这种语义，可以减少 \code{+1} 和 \code{-1} 在业务代码里到处出现。

线段树把数组按区间递归划分成二叉树。递归概念容易理解，但 C++ 实现可以用迭代线段树避免调用栈。叶子放在底层，父节点保存两个孩子的合并结果。点更新时从叶子改起一路向上重算；区间查询时左右边界向上跳，收集覆盖查询区间的若干节点。它适合动态区间最值、区间和、扫描线、最长连续段、括号匹配信息等。

\section{图表示：邻接矩阵、邻接表、边表和状态图}

图结构的核心是节点和边。邻接矩阵适合点数小且边很多的图，判断两点是否相邻是常数，但空间是 \complexity{O(n^2)}。邻接表适合稀疏图，空间 \complexity{O(n+m)}，遍历某点邻居高效。边表适合 Kruskal、Bellman-Ford 和离线处理，因为算法按边扫描而不是按点找邻居。网格图可以不显式建图，用方向数组现场生成邻居。

建图题最容易错在边方向和状态维度。课程表中，如果边从课程指向先修课，拓扑输出顺序就和学习顺序相反；单词接龙中，如果每次找邻居都扫描所有单词，BFS 外层正确也会超时；障碍消除最短路中，状态不是位置，而是位置加剩余消除次数；访问所有节点的最短路径中，状态是当前位置加已访问集合。状态图的节点可能不是输入中的一个对象，而是多个变量的组合。

图结构在实际系统中非常常见。编译器构建控制流图和调用图，包管理器维护依赖图，数据库查询优化器处理算子图，网络路由维护拓扑图，操作系统调度和资源等待可以形成等待图。刷题中的 BFS、拓扑排序和最短路，是这些系统问题的简化模型。

\section{并查集：等价关系、扩展权值和回滚}

并查集维护的是等价关系。每个集合有一个代表元，\code{find} 返回代表，\code{unite} 把两个代表连接起来。路径压缩降低未来查找成本，按大小合并避免树退化。它适合动态合并、判断同组、统计连通块、检测无向环和 Kruskal 最小生成树。

并查集的变体很多。带权并查集维护节点到父节点的权值关系，可用于除法求值、相对距离和势能差；奇偶并查集维护节点到根的奇偶关系，可用于二分图约束和敌友关系；带大小并查集维护集合规模；回滚并查集支持撤销合并，常用于离线动态连通性。每个变体都遵循同一原则：合并代表时，额外信息必须同步更新到新的父子关系上。

\begin{lstlisting}[language=C++,caption={并查集统计连通块数量}]
class ComponentSet {
public:
    explicit ComponentSet(int n)
        : parent_(static_cast<std::size_t>(n)),
          size_(static_cast<std::size_t>(n), 1),
          components_(n)
    {
        std::iota(this->parent_.begin(), this->parent_.end(), 0);
    }

    bool unite(int lhs, int rhs)
    {
        int root_lhs = this->find(lhs);
        int root_rhs = this->find(rhs);
        if (root_lhs == root_rhs) {
            return false;
        }
        if (this->size_[static_cast<std::size_t>(root_lhs)] <
            this->size_[static_cast<std::size_t>(root_rhs)]) {
            std::swap(root_lhs, root_rhs);
        }
        this->parent_[static_cast<std::size_t>(root_rhs)] = root_lhs;
        this->size_[static_cast<std::size_t>(root_lhs)] +=
            this->size_[static_cast<std::size_t>(root_rhs)];
        --this->components_;
        return true;
    }

    int components() const { return this->components_; }

private:
    int find(int x)
    {
        int root = x;
        while (this->parent_[static_cast<std::size_t>(root)] != root) {
            root = this->parent_[static_cast<std::size_t>(root)];
        }
        while (this->parent_[static_cast<std::size_t>(x)] != x) {
            const int next = this->parent_[static_cast<std::size_t>(x)];
            this->parent_[static_cast<std::size_t>(x)] = root;
            x = next;
        }
        return root;
    }

    std::vector<int> parent_;
    std::vector<int> size_;
    int components_ = 0;
};
\end{lstlisting}

这段代码把连通块数量作为并查集状态维护。每次成功合并两个不同集合，连通块数减一；重复合并不改变数量。省份数量、网络连通操作次数、岛屿动态添加都可以使用这个模式。若题目需要知道集合内最大值、边数或权值差，只需要在根节点上维护额外数组，并在合并时更新。

\section{非线性结构的应用谱和实验}

哈希表应用：两数之和、字母异位词分组、前缀和计数、克隆图、LRU 索引、字符串到编号映射。堆应用：Top K、动态中位数、Dijkstra、多路归并、任务调度、最小区间。平衡树应用：日程安排、前驱后继、动态区间、滑动窗口有序统计。Trie 应用：前缀搜索、单词搜索 II、最大异或、敏感词过滤。树状数组和线段树应用：逆序对、区间和、扫描线、动态排名、范围更新。图应用：依赖分析、最短路、连通块、拓扑排序。并查集应用：动态连通、冗余连接、账户合并、等式约束、最小生成树。

\begin{labbox}
非线性结构实验：第一，给 \code{unordered_map} 设置不同 \code{reserve}，统计 rehash 次数和运行时间。第二，用完整排序、大小为 K 的堆、快速选择分别求第 K 大，比较不同 K 下的表现。第三，用 \code{map} 实现日程安排，再用无序表尝试，解释为什么无序表无法找相邻区间。第四，用树状数组统计逆序对，并手算离散化后的更新和查询。第五，把并查集的路径压缩关掉，构造链式合并，观察查找路径长度变化。
\end{labbox}

\begin{principlebox}{用 LeetCode 案例选择非线性结构}
LeetCode 560 和为 K 的子数组需要历史前缀精确查询，用哈希表；LeetCode 215 数组第 K 大和 LeetCode 295 数据流中位数需要动态极值或中位边界，用堆；LeetCode 731 我的日程安排 II 需要按端点顺序扫描，用有序表或线段树；LeetCode 212 单词搜索 II 需要大量前缀共享，用 Trie；LeetCode 307 区域和检索需要动态前缀聚合，用树状数组或线段树；LeetCode 200 岛屿数量需要可达性，用图搜索；LeetCode 721 账户合并只需要合并集合和判断同组，用并查集。每次选择结构都要同时说明它不能做什么，避免把结构能力无限放大。
\end{principlebox}

\section{非线性结构变种题谱}

哈希表变种可以按 value 语义分类。保存下标：两数之和、同构字符串、最长无重复子串的上次出现位置。保存次数：异位词分组、前缀和计数、优美子数组、频率栈。保存最早位置：最长和为 K、连续子数组和、最长等量零一子数组。保存集合代表：并查集邮箱映射、字符串编号、图克隆映射。保存对象迭代器：LRU 缓存、全 O(1) 数据结构。题解里必须写清 value 语义，否则变体一来就会把次数、位置和布尔存在混用。

堆变种可以按“极值来源”分类。固定大小堆维护 Top K，双堆维护中位数，多路堆维护 K 个有序流的当前头，带过期堆处理滑动窗口和 Dijkstra 旧状态，贪心堆在任务调度和 IPO 中选择当前可执行最优项目。每类堆都要补一个过期或比较器问题：Top K 要说明堆顶代表边界；双堆要说明大小和平衡；多路堆要说明堆元素里保存来自哪一路；带过期堆要说明过期检查发生在使用堆顶之前。

平衡树和有序结构变种按查询目的分类。前驱后继类：存在重复元素 III、时间键值存储、最近座位。区间维护类：我的日程安排、插入区间、合并数据流区间。排名统计类：逆序对、区间和个数、滑动窗口中位数。扫描线类：会议室、天际线、矩形面积。若只需要存在性，哈希表更简单；一旦需要“最近的更大/更小”“第 k 个”“区间相邻关系”，就要转向有序结构、树状数组或线段树。

Trie 变种按边的含义分类。字符 Trie 用于单词前缀和字典匹配；反向 Trie 用于后缀查询和流式字符匹配；压缩 Trie 用于减少长单链空间；二进制 Trie 用于最大异或、按位贪心和网络前缀匹配；AC 自动机把 Trie 和失败指针结合，用于多模式匹配。刷题中最常见的是字符 Trie 和二进制 Trie。二进制 Trie 的边是 0/1 位，不是字符；每次为了最大异或，优先走当前位相反的分支。

树状数组和线段树变种按更新和查询分类。点更新加区间查询是最基础形态；区间更新加点查询可以用差分思想；区间更新加区间查询需要懒标记或两个树状数组；区间最大值、最小值、最大子段和、最长连续一都需要自定义合并信息。线段树难点不在模板长，而在节点信息能否合并。只要父节点信息可以由左右孩子合成，就可能使用线段树；如果合并依赖更远上下文，就需要换建模。

图变种按边权和状态分类。无权图最短路用 BFS，零一权图用 0-1 BFS，非负权图用 Dijkstra，有负权边但无负环可用 Bellman-Ford 或 SPFA 思想，有向无环图可用拓扑 DP。状态图题要明确状态维度：开锁是四位字符串，访问所有节点是节点加集合，障碍消除是位置加剩余次数，带钥匙迷宫是位置加钥匙掩码。状态维度决定 visited 数组维度，也决定复杂度。

并查集变种按附加信息分类。普通并查集维护连通块，带大小并查集维护集合规模，带权并查集维护相对比值或距离，奇偶并查集维护二分关系，网格并查集把二维坐标映射成一维编号，回滚并查集支持撤销。每个变种都要回答：路径压缩是否影响附加信息，合并两个根时附加信息如何更新，查询需要的是根信息还是节点到根的信息。

\section{工程应用：从刷题结构到系统结构}

哈希表和链表组合构成缓存系统的基本骨架。LRU 用哈希表定位节点，用双向链表维护新旧顺序；LFU 在 LRU 基础上增加频次维度，用频次到链表的映射维护同频内部顺序。这个结构也解释了为什么单独一个哈希表或单独一个链表都不够：哈希表不知道谁最旧，链表不能常数定位 key。

堆和有序树常用于调度。操作系统调度器、定时器、任务队列、网络包优先级，都需要从动态集合中取出某种最优对象。简单定时任务可以用小顶堆按到期时间排序；需要频繁删除和修改优先级时，可能要用有序树或索引堆；需要公平性时，键可能不是绝对时间，而是虚拟运行时间。刷题中的任务调度器和会议室问题，是这些系统问题的简化模型。

Trie 和哈希索引用于检索系统。搜索建议、自动补全、敏感词过滤、路由前缀匹配，都依赖前缀共享或快速定位。普通哈希表适合完整 key 查询；Trie 适合前缀查询；倒排索引适合从词到文档集合；布隆过滤器适合快速判断“肯定不存在或可能存在”。刷题不一定直接考布隆过滤器，但理解这些结构能扩展你对“查询模型”的认识。

图结构用于依赖和流动。包管理依赖、编译依赖、课程安排、任务 DAG、网络路由、社交关系和知识图谱都可以建成图。拓扑排序回答依赖顺序，最短路回答最小代价，连通块回答分组，可达性回答影响范围。真正做工程时，图的节点和边往往来自业务对象，需要先把业务关系抽象成图，再选算法。

\begin{principlebox}{变种题迁移标准}
看到变种题时，不要先问“它像哪道 LeetCode 题”，先问四个问题：查询对象是什么，更新对象是什么，是否需要顺序，是否需要合并历史状态。答案分别会指向哈希、堆、有序结构、Trie、区间树、图或并查集。能从访问模式推出结构，才算真正掌握数据结构。
\end{principlebox}

\section{哈希结构的深层语义：键、值和状态等价}

哈希表题真正难的地方不是调用 \code{find}，而是定义 key 和 value。key 必须表示“哪些状态在题目里应当被看作同一类”，value 必须表示“遇到同一类状态时还需要保留什么信息”。两数之和的 key 是数值，value 是下标；字母异位词分组的 key 是字符频次签名，value 是字符串列表；同构字符串的 key 可以是字符映射，value 是另一侧字符；子数组和的 key 是前缀和，value 可能是次数、最早位置或最近位置。key 定义错，哈希表越快越会快地得到错误答案。

状态等价尤其重要。字母异位词里，\code{"eat"} 和 \code{"tea"} 的原字符串不同，但字符频次相同，所以它们在分组问题中等价；快乐数里，一个整数经过平方和转移后如果再次出现，就说明进入循环，所以 key 是当前数值状态；克隆图里，原节点地址决定唯一身份，节点值可能重复，所以 key 应该是原节点指针或稳定编号，而不是节点值。经典教材讲哈希时会强调抽象等价关系，因为这比容器 API 更接近算法本质。

字母异位词分组先从暴力说起。最直接的做法是拿每个单词和已有分组逐个比较：若两个单词排序后相同，或逐字符计数后相同，就放进同一组。这个版本的瓶颈有两层：第一，单词之间会反复两两比较；第二，同一个单词的规范表示可能被重复计算。优化的关键是把“是否属于同一组”改写成一个状态等价 key：异位词共享同一份字符频次签名。每个单词只计算一次签名，然后用哈希表或有序 map 直接进入对应桶。这样，比较单词对的问题就变成了计算 key 并追加到 value 列表。

\begin{principlebox}{图示：异位词分组的状态等价}
\begin{tabular}{@{}p{0.2\linewidth}p{0.72\linewidth}@{}}
原始单词 & \code{eat}、\code{tea}、\code{ate} 是不同字符串。\\
规范签名 & 三者字符频次都为 \code{a:1, e:1, t:1}。\\
哈希键 & 频次签名作为 key，表示“属于同一组”的等价状态。\\
哈希值 & value 保存这一组所有原始字符串，输出时直接取桶。\\
\end{tabular}
\end{principlebox}

\begin{lstlisting}[language=C++,caption={异位词分组：频次签名定义状态等价}]
std::vector<std::vector<std::string>> group_anagrams(
    const std::vector<std::string>& words)
{
    constexpr int ALPHABET_SIZE = 26;
    std::map<std::array<int, ALPHABET_SIZE>, std::vector<std::string>> groups;

    for (const std::string& word : words) {
        std::array<int, ALPHABET_SIZE> signature{};
        for (const char ch : word) {
            ++signature[ch - 'a'];
        }
        groups[signature].push_back(word);
    }

    std::vector<std::vector<std::string>> answer;
    answer.reserve(groups.size());
    for (auto& entry : groups) {
        auto& group = entry.second;
        answer.push_back(std::move(group));
    }
    return answer;
}
\end{lstlisting}

这里用 \code{std::map} 是为了避免给 \code{array} 写自定义哈希，同时让输出顺序稳定；若追求平均性能，可以换成 \code{unordered_map} 并提供哈希函数。这个例子还说明 value 的类型由题目输出决定：如果只问有多少组，value 可以是计数；如果要返回每组字符串，value 必须保存列表。不要把“哈希表”当成固定模板，它只是实现 key 到 value 的映射。

哈希表的工程问题也不能省略。第一，扩容会 rehash，迭代器失效，复杂对象的移动和分配成本会上升。第二，自定义 key 要保证相等对象具有相同哈希值，否则容器行为不符合预期。第三，浮点数不适合作为普通哈希 key，精度误差会让数学上相等的状态变成不同 key。第四，如果输入可能被恶意构造，哈希退化会影响性能。刷题环境通常不要求处理所有工程问题，但写书应该把边界讲出来。

\section{堆结构的边界：局部有序不是全局有序}

堆的最重要边界是：它只保证堆顶最优，不保证内部整体有序。很多学生误以为从 \code{priority_queue} 里遍历出来就是有序序列，这是错的。堆适合反复取当前极值，不适合查找任意元素，不适合直接删除任意元素，不适合回答前驱后继。如果题目需要“窗口内第 k 小”“动态排名”“相邻区间”，堆往往不够，需要有序结构、双堆加懒删除，或树状数组和线段树。

双堆维护中位数是堆边界的经典扩展。一个大顶堆保存较小一半，一个小顶堆保存较大一半。中位数来自两个堆顶，而不是某一个全局有序数组。不变量有两个：大小不变量和顺序不变量。大小不变量保证两个堆元素数量差不超过一；顺序不变量保证左堆每个元素都不大于右堆每个元素。只维护大小会出错，因为堆顶可能交叉。

数据流中位数的暴力版本很自然：每来一个数就插入数组，查询中位数时排序或线性选择。若每次完整排序，旧元素之间已经确定的相对关系被反复计算；若维护一个有序数组，插入位置可以二分找到，但中间插入仍要搬移大量元素。中位数只需要中间边界，而不是完整有序序列。于是把数据切成两半：较小一半只关心最大值，较大一半只关心最小值。左半用大顶堆，右半用小顶堆，插入后重平衡，就把“维护全序”降成“维护两侧边界”。

\begin{principlebox}{图示：数据流中位数的双堆边界}
\begin{tabular}{@{}p{0.2\linewidth}p{0.72\linewidth}@{}}
左半部分 & 大顶堆 \code{lower} 保存较小一半，堆顶是左半最大值。\\
右半部分 & 小顶堆 \code{upper} 保存较大一半，堆顶是右半最小值。\\
顺序不变量 & \code{lower.top() <= upper.top()}，左半所有元素不大于右半。\\
大小不变量 & 两堆大小差不超过一；本书示例让左堆可以多一个。\\
中位数 & 奇数取左堆顶，偶数取两个堆顶平均值。\\
\end{tabular}
\end{principlebox}

\begin{lstlisting}[language=C++,caption={数据流中位数：双堆维护大小和顺序}]
class MedianFinder {
public:
    void add_num(int value)
    {
        if (this->lower_.empty() || value <= this->lower_.top()) {
            this->lower_.push(value);
        } else {
            this->upper_.push(value);
        }
        this->rebalance();
    }

    double find_median() const
    {
        if (this->lower_.size() == this->upper_.size()) {
            return (static_cast<double>(this->lower_.top()) +
                    static_cast<double>(this->upper_.top())) /
                   2.0;
        }
        return static_cast<double>(this->lower_.top());
    }

private:
    void rebalance()
    {
        if (this->lower_.size() < this->upper_.size()) {
            this->lower_.push(this->upper_.top());
            this->upper_.pop();
        } else if (this->lower_.size() > this->upper_.size() + 1) {
            this->upper_.push(this->lower_.top());
            this->lower_.pop();
        }
    }

    std::priority_queue<int> lower_;
    std::priority_queue<int, std::vector<int>, std::greater<int>> upper_;
};
\end{lstlisting}

这段代码规定左堆元素数可以比右堆多一个，因此奇数个元素时中位数在左堆顶。插入时先按值进入对应堆，再平衡大小。若题目要求删除窗口外元素，就要引入懒删除表，并在访问堆顶前清理过期元素。双堆题的难点不是写堆，而是维护不变量并在删除时保持两个堆的有效大小。

堆在工程里常用于调度和归并。定时器系统可以把任务按到期时间放入小顶堆；日志合并可以把多个有序文件的当前行放入小顶堆；Dijkstra 把候选节点按当前距离放入小顶堆；操作系统或线程池可以用优先级队列选择下一项任务。所有这些场景都有同一结构：候选集合动态变化，但每次只需要当前最优候选。若还需要修改任意候选的优先级，普通堆就不够，要用索引堆、有序树或重新插入加过期检查。

\section{有序结构和区间结构：从相邻关系到聚合信息}

有序结构回答的是“在顺序上的邻居是谁”。\code{lower_bound} 找到第一个不小于目标的元素，配合前一个迭代器，就能检查插入区间是否和相邻区间冲突。会议室、日程安排、数据流区间、最近座位、存在重复元素 III，都依赖相邻关系。哈希表再快，也不能告诉你一个值在有序集合中的前驱和后继。

区间结构回答的是“一个范围内的聚合值是什么”。树状数组擅长前缀可逆聚合，线段树擅长更通用的区间合并。选择它们时先问两个问题：更新是点还是区间，查询是点还是区间。点更新加区间查询，树状数组简单；区间更新加区间查询，线段树懒标记更直接；如果只做静态查询，前缀和或稀疏表可能更合适。不要一看到区间就写线段树，很多题只需要排序加扫描线或前缀和。

区间最大值如果数组不变，暴力每次扫描区间就能回答，只是多次查询会重复访问同一段；前缀和只能处理可逆的和，不能直接由两个前缀最大值相减得到区间最大。若又有点更新，静态稀疏表也不合适。线段树的推导来自“区间答案能由左右子区间合并”：一个区间的最大值等于左右半区最大值的较大者。把这个合并关系按层缓存起来，点更新只影响从叶子到根的一条路径，区间查询只需收集若干完整覆盖的小区间。

\begin{principlebox}{图示：线段树的区间合并}
\begin{tabular}{@{}p{0.2\linewidth}p{0.72\linewidth}@{}}
叶子节点 & 保存单个数组位置的值。\\
父节点 & 保存左右孩子区间的合并结果，例如 \code{max(left, right)}。\\
点更新 & 修改一个叶子，只沿父指针向上重算一条路径。\\
区间查询 & 把查询区间拆成若干不相交的完整节点区间，再合并这些节点。\\
成立前提 & 父区间答案必须能由左右子区间答案合成。\\
\end{tabular}
\end{principlebox}

\begin{lstlisting}[language=C++,caption={迭代线段树：点更新和区间最大值}]
class SegmentTree {
public:
    explicit SegmentTree(const std::vector<int>& values)
        : size_(static_cast<int>(values.size())),
          tree_(static_cast<std::size_t>(values.size() * 2), 0)
    {
        for (int i = 0; i < this->size_; ++i) {
            this->tree_[static_cast<std::size_t>(this->size_ + i)] = values[i];
        }
        for (int i = this->size_ - 1; i > 0; --i) {
            this->tree_[static_cast<std::size_t>(i)] =
                std::max(this->tree_[static_cast<std::size_t>(i * 2)],
                         this->tree_[static_cast<std::size_t>(i * 2 + 1)]);
        }
    }

    void update(int index, int value)
    {
        int position = index + this->size_;
        this->tree_[static_cast<std::size_t>(position)] = value;
        for (position /= 2; position > 0; position /= 2) {
            this->tree_[static_cast<std::size_t>(position)] =
                std::max(this->tree_[static_cast<std::size_t>(position * 2)],
                         this->tree_[static_cast<std::size_t>(position * 2 + 1)]);
        }
    }

    int range_max(int left, int right) const
    {
        constexpr int NEGATIVE_INF = std::numeric_limits<int>::min();
        int answer = NEGATIVE_INF;
        int l = left + this->size_;
        int r = right + this->size_;
        while (l < r) {
            if ((l & 1) != 0) {
                answer = std::max(answer, this->tree_[static_cast<std::size_t>(l)]);
                ++l;
            }
            if ((r & 1) != 0) {
                --r;
                answer = std::max(answer, this->tree_[static_cast<std::size_t>(r)]);
            }
            l /= 2;
            r /= 2;
        }
        return answer;
    }

private:
    int size_ = 0;
    std::vector<int> tree_;
};
\end{lstlisting}

这个实现使用半开区间 \code{[left, right)}。查询时，若左边界是右孩子，就把它纳入答案并右移；若右边界是右孩子的后一位，就先左移再纳入答案。迭代线段树比递归版本短，也避免 C++ 调用栈问题，但它要求你非常清楚半开区间语义。若题目需要区间加法和区间最大值，就要增加懒标记；若节点信息不是一个整数，而是最大子段和、前缀最大、后缀最大和总和，就要定义可合并的结构体。

\section{树结构：从遍历顺序到信息流}

树题不应该只按前序、中序、后序背模板，而要按信息流分类。自顶向下题把父节点信息传给子节点，例如路径和、验证 BST 上下界、二叉树所有路径；自底向上题把子树信息汇总给父节点，例如高度、直径、平衡二叉树、最大路径和；层次题按距离或深度组织节点，例如层序遍历、右视图、最小深度；结构改造题改变节点链接，例如展开二叉树、反转、删除 BST 节点。遍历顺序服务于信息流，而不是目的本身。

二叉搜索树的中序递增只是它的一种表象。更本质的性质是：每个节点把值域划分成左右两个区间。验证 BST 可以中序，也可以传上下界；第 K 小可以中序计数；删除节点时，如果左右孩子都存在，可以用后继节点替换当前值，再删除后继节点；最近公共祖先在 BST 中可以利用值域方向，而普通二叉树必须在左右子树中寻找目标。把值域约束讲清楚，很多 BST 题就能串起来。

树形 DP 的核心是定义返回值。最大路径和中，返回给父节点的不是当前子树内部最大路径，而是“从当前节点向下延伸到某一侧的最大贡献”；全局答案才记录可能经过当前节点同时连接左右两侧的路径。打家劫舍 III 中，一个节点的状态至少要区分偷和不偷；若只返回一个最大值，会丢失父子不能同时选的约束。树题最常见的错误就是返回值语义太含糊。

\section{图结构：状态空间、边权和访问时机}

图题的第一步是定义状态，而不是选择 DFS 或 BFS。状态可能是一个节点，也可能是多个变量的组合。普通岛屿题的状态是坐标；障碍消除最短路的状态是坐标加剩余消除次数；访问所有节点的最短路径是当前节点加访问集合；带钥匙迷宫是坐标加钥匙掩码；单词接龙是一个字符串或字符串编号。状态定义少了维度，会把不同情况错误合并；状态定义多了无用维度，会让复杂度爆炸。

访问标记的时机也很关键。普通 BFS 中，通常在入队时标记已访问，而不是出队时标记。因为一个节点可能被多个邻居发现，如果等出队才标记，它会被重复入队。Dijkstra 则不同，一个节点可以用更短距离再次进入堆，弹出时检查当前距离是否过期。并查集又不同，它没有访问队列，而是用集合代表压缩连通关系。不同结构有不同“状态已经确定”的时刻，不能混用。

障碍消除最短路不能直接把状态写成坐标。暴力搜索会枚举从起点出发的所有路径，并在路径里记录已经消除过多少障碍；普通 BFS 若只用 \code{visited[row][col]}，会把“到达同一格子但剩余消除次数不同”的两种情况错误合并。瓶颈不是 BFS 本身，而是状态定义不完整导致要么漏解，要么保留太多等价路径。优化后的状态至少包含位置和剩余资源；进一步观察，若两条路径到达同一格子，距离按 BFS 层已经不更差，而剩余消除次数更多的状态支配剩余更少的状态。因此每个格子只保存见过的最大剩余次数，就能剪掉未来不可能更好的状态。

\begin{principlebox}{图示：状态图 BFS 的支配关系}
\begin{tabular}{@{}p{0.2\linewidth}p{0.72\linewidth}@{}}
错误状态 & 只保存 \code{(row, col)}，会把资源不同的路径合并。\\
完整状态 & 保存 \code{(row, col, remaining)}，位置和剩余消除次数共同决定未来能力。\\
支配剪枝 & 同一格子上，剩余次数更少且距离不更短的状态没有保留价值。\\
压缩记录 & \code{best_remaining[row][col]} 保存到达该格子时见过的最大剩余资源。\\
BFS 语义 & 队列按步数扩展，第一次到达终点时步数最短。\\
\end{tabular}
\end{principlebox}

\begin{lstlisting}[language=C++,caption={状态图 BFS：位置加剩余资源才是完整状态}]
int shortest_path_with_elimination(
    const std::vector<std::vector<int>>& grid,
    int eliminations)
{
    if (grid.empty() || grid[0].empty()) {
        return -1;
    }

    constexpr std::array<std::pair<int, int>, 4> DIRECTIONS{
        std::pair{1, 0}, std::pair{-1, 0},
        std::pair{0, 1}, std::pair{0, -1}
    };

    const int rows = static_cast<int>(grid.size());
    const int cols = static_cast<int>(grid[0].size());
    std::vector<std::vector<int>> best_remaining(
        static_cast<std::size_t>(rows),
        std::vector<int>(static_cast<std::size_t>(cols), -1));

    std::queue<std::array<int, 4>> queue;
    queue.push({0, 0, eliminations, 0});
    best_remaining[0][0] = eliminations;

    while (!queue.empty()) {
        const auto [row, col, remaining, distance] = queue.front();
        queue.pop();
        if (row == rows - 1 && col == cols - 1) {
            return distance;
        }

        for (const auto& [dr, dc] : DIRECTIONS) {
            const int nr = row + dr;
            const int nc = col + dc;
            if (nr < 0 || nr >= rows || nc < 0 || nc >= cols) {
                continue;
            }
            const int next_remaining =
                remaining - grid[static_cast<std::size_t>(nr)]
                                  [static_cast<std::size_t>(nc)];
            if (next_remaining < 0 ||
                next_remaining <= best_remaining[static_cast<std::size_t>(nr)]
                                                 [static_cast<std::size_t>(nc)]) {
                continue;
            }
            best_remaining[static_cast<std::size_t>(nr)]
                          [static_cast<std::size_t>(nc)] = next_remaining;
            queue.push({nr, nc, next_remaining, distance + 1});
        }
    }
    return -1;
}
\end{lstlisting}

这里没有用三维 \code{visited[row][col][remaining]}，而是为每个格子保存到达它时剩余消除次数的最好值。若再次到达同一格子但剩余次数更少或相等，未来不会比之前更好，可以剪掉。这个剪枝来自状态支配关系：位置相同、距离不更短、资源更少的状态没有保留价值。状态图题经常有这种优化，但前提是你能证明一个状态支配另一个状态。

边权决定最短路算法。无权边用 BFS；0/1 权边用 0-1 BFS；非负权边用 Dijkstra；存在负权边时，Dijkstra 的贪心确定性失效，需要 Bellman-Ford 或其他算法；DAG 最短路可以按拓扑序 DP。题目如果写的是“每一步代价相同”，才是普通 BFS；如果每条边代价不同，即使看起来也是从一个点走到另一个点，也不能直接用普通队列。

\section{并查集：等价类、势能和离线思维}

并查集维护的是等价类。它能快速回答两个元素是否已经连通，却不能告诉你连通路径是什么，也不能处理删除边后的在线连通性。这个能力边界很重要。省份数量、冗余连接、账户合并、等式可满足性，都只需要合并和查询同组；最短路径、路径条数、删除某条边后是否仍连通，就不是普通并查集擅长的范围。

带权并查集把“到父节点的关系”也保存下来。除法求值中，若 \code{a / b = 2}，可以把 a 和 b 放到同一集合，并维护每个节点到根的比例。查询 \code{x / y} 时，如果两者根不同，答案不存在；根相同，则用 \code{x/root} 除以 \code{y/root}。奇偶并查集也是同理，只是权值从比例变成奇偶关系。难点在合并两个根时，要根据已知约束推导新父子关系上的权值。

回滚并查集和离线动态连通则体现了并查集的高级用法。普通路径压缩会破坏历史，不适合回滚；如果要支持撤销，常用按大小合并并记录被修改的父亲和大小，撤销时恢复。区间时间上的边添加和删除可以用线段树分治把边分配到生效时间段，再配合回滚并查集处理查询。面试不常要求手写完整实现，但理解这个方向能帮助你知道普通并查集为什么不能删除边，以及离线算法如何绕开在线限制。

\section{数据结构组合：一道题往往不只一种结构}

很多高质量面试题不是考单个结构，而是考组合。LRU 是哈希表加双向链表；LFU 是哈希表加频次表加链表；Twitter 设计题是哈希表加链表或数组加堆归并；单词搜索 II 是 Trie 加 DFS 回溯；账户合并是哈希表加并查集；滑动窗口中位数是双堆加延迟删除哈希表；最小覆盖区间是排序加堆或多指针；天际线是扫描线加堆或有序表。组合题要求你给每个结构分配明确职责。

设计题可以按数据流拆解。先列出操作：插入、删除、查询、移动、合并、取极值、取前驱后继、按前缀查找。再为每个操作找结构能力。若两个操作互相冲突，就用两个结构互相索引，并维护一致性。LRU 中，哈希表和链表必须同时更新；Twitter 中，用户到推文列表的映射和关注关系必须分开；滑动窗口中位数中，堆里有旧值，延迟删除表必须在取堆顶前清理。组合题的 bug 往往来自一个结构更新了，另一个结构没同步。

\begin{principlebox}{组合结构的写法}
先写每个结构的职责，再写同步规则。职责包括：谁负责定位，谁负责顺序，谁负责极值，谁负责连通性，谁负责前缀，谁负责区间。同步规则包括：插入时更新哪些结构，删除时更新哪些结构，移动时是否保持迭代器有效，查询前是否要清理过期状态。没有同步规则的组合结构很容易在边界输入上失效。
\end{principlebox}

\section{非线性结构题目的讲解标准}

非线性结构题讲解至少要回答六个问题。第一，状态或对象是什么。第二，结构里的 key、节点、边、堆元素或集合代表分别表示什么。第三，暴力方法慢在哪里，是查找历史、找极值、找相邻、查前缀、查区间，还是判断连通。第四，所选结构把哪一步降到了什么复杂度。第五，结构的限制是什么，比如哈希无序、堆不能任意删除、并查集不能给路径。第六，变体来临时如何调整 value、比较器、状态维度或附加信息。

以课程表为例，对象是课程，边是先修关系，暴力慢在反复找没有依赖的课程；邻接表保存从一门课指向依赖它的后续课程，入度数组保存未满足依赖数量，队列保存当前可学课程。每弹出一门课，就减少后续课程入度。若最终弹出数量不足，说明有环。变体课程表 II 需要输出顺序，只需记录弹出序列；若课程有持续时间并要求最短完成时间，就要在拓扑序上做 DP；若依赖会动态增删，普通拓扑排序就不够。

以账户合并为例，对象是账户和邮箱。哈希表把邮箱映射到第一次出现的账户，并查集合并共享邮箱的账户，最后按根收集邮箱。暴力方法可能两两比较账户邮箱集合，复杂度很高；哈希表让共享邮箱检测变成平均常数，并查集让账户连通分量合并接近常数。变体如果要求按用户名区分，key 就必须包含用户名或先按用户名分组；如果邮箱大小写不敏感，就要先标准化邮箱。题解必须把这些建模前提说清楚。

以单词搜索 II 为例，对象是棋盘路径和单词集合。暴力做法是对每个单词在棋盘上搜索一次，重复遍历大量相同前缀。Trie 把所有单词的公共前缀合并，DFS 当前路径同时对应 Trie 中一个节点；如果当前字符没有对应边，就立即剪枝。这里 Trie 负责前缀查询，DFS 负责空间路径，访问标记负责同一单元格不能重复使用。变体如果允许八方向移动，只改邻居生成；如果允许重复使用格子，访问标记规则要改；如果字符集很大，Trie 孩子结构也要改。

把这些标准坚持下来，非线性结构就不会变成孤立模板。哈希、堆、树、图、并查集只是不同查询模型的工具；教材的任务是让读者知道每个工具解决什么瓶颈、付出什么代价、在什么条件下失效，以及如何和其他结构组合。

